\section{Background and Challenges}
\label{sec:motivation}

\textbf{Mixed-precision weights on a bit-serial datapath.}
FlexPosit~\cite{gao2026flexposit} applies posit formats~\cite{gustafson2017posit} to LLM weights and lets an accelerator choose a precision from P4 to P8 per group of output channels.
For Qwen2.5-14B~\cite{qwen2025qwen25}, uniform P4 and P5 weights raise WikiText-2 perplexity from 5.88 (FP16) to 7.86 and 6.41, while a Fisher-guided allocation of P4--P8 per 32-channel group reaches 6.11 at 5.5 bits per weight; the key and value projections average 7.2 and 7.6 bits, most others P5 or P6.
Serializing weight bits costs decode no throughput, since decode is HBM-bound at every precision (\Cref{fig:motivation}b), and makes MAC time and energy scale with precision. Each array command of our datapath processes one 32-column weight tile and serializes its weight bits, so a tile takes $94+20(p-4)$ cycles at precision $p$, its MAC energy grows from 1.24 to 2.04\,pJ per MAC, and its weight bytes grow in proportion to $p$.
We call one node of the operator graph, such as the query projection of one layer, a \emph{block}. A block's compute time $C$ therefore follows the precisions of its tiles and its HBM time $M$ its weight bytes.
DVFS can lower a block's operating point without delay as long as its stretched $C$ stays within $M$; we call the lowest such point its \emph{free point}, and the latency budget under the TTFT and TPOT objectives the \emph{slack}.

\begin{figure}[t]
\centering
\includegraphics[width=\linewidth]{figures/fig_motivation.pdf}
\caption{(a) Precision-dependent imbalance: saving over power and clock gating of a single rail switched per phase or per layer, and of \sys{}'s multi-domain plan (calibrated first-order model, 3$\times$2K-token requests). (b) Balance shift: compute time over HBM time of Qwen2.5-14B projections on 256 arrays; precision moves the prefill query projection across the balance point.}
\label{fig:motivation}
\end{figure}

\begin{figure}[t]
\centering
\includegraphics[width=\linewidth]{figures/fig_units.pdf}
\caption{Precision moves each domain differently (Qwen2.5-14B, 2K-token prompt, 256 arrays, peak V/$f$; Mix is the 5.5-bit Fisher allocation): (a) prefill busy time relative to P4; (b) energy under power and clock gating.}
\label{fig:units}
\end{figure}

\textbf{C1: Precision-dependent imbalance.}
We ran 45 nominal cases, three models at three prompt lengths and five precisions (\Cref{fig:units}).
From P4 to P8 a Qwen2.5-14B prefill's HBM time grows by 94\%, MAC time by 33\%, memory-logic time by 11\%, and vector time not at all; the nine model and prompt combinations agree (HBM +91--95\%, MAC +33--56\%, memory logic +11\%, vector 0).
$C/M$ falls from 1.16 to 0.80 and the compute-bound share of prefill blocks from 0.74 to 0 on Qwen2.5-14B.
Energy moves with it, from memory logic (32 to 25\%) to MAC (35 to 41\%) and HBM (19 to 23\%), and the vector rail, which dominates prefill busy time, keeps TTFT from moving as much as the blocks do (+20\%).
The components therefore want different voltages at once: at a 5\% TTFT budget the planner puts the vector rail near 0.69\,V at every precision, the MAC rail at 0.55--0.59\,V and LOGIC near 0.47\,V, and at 25\% the MAC rail reaches the 0.45\,V floor while the vector rail stays at 0.69\,V.
The naive solution is one rail per core, as in per-core NPU DVFS~\cite{ascend_dvfs:asplos25}.
That rail must stay high for the vector unit and so cannot lower the MAC domain.
Even with 34 points and full knowledge of the requests, switched per phase it saves 1.2--3.9\% at TTFT budgets of 2--25\%, and switched per layer 0.8--1.7\% (\Cref{fig:motivation}a).


\textbf{C2: Crossing latency.}
Solving C1 with separate domains creates the next problem. Precision shortens matrix tiles but not the vector chains between them, and the domains insert synchronizers on these chains: a dependent operation that waits for a result from another domain pays the crossing twice, about 29\,ns per instruction on a SiLU chain against a 4.2\,ns local issue in the FP32 vector unit (VU).
Normalization, activation and softmax kernels are chains of dependent vector instructions, so these round trips add up: with every dependency crossing, the kernels take 4.2--6.9$\times$ as long and TTFT grows by 41--46\% across nine workloads.
The naive fix keeps a shadow copy of recent results in the VU domain but aliases it to frontend registers~\cite{xue2026finedvfs}, so every result is still written back across the crossing; the return queue must cover the round trip of every instruction.

\textbf{C3: Balance shift.}
Precision decides on which side of the balance point ($C=M$) a block falls (\Cref{fig:motivation}b).
Decode blocks read every weight once per token and stay HBM-bound at any precision ($C/M$ 0.32--0.42); prefill blocks of a 1920-token prompt sit near the balance point, and the query projection is compute-bound from P4 to P6 but HBM-bound at P8.
Prefill holds 93.8\% of the MAC energy of a 1920+128-token request, so this side decides most of what DVFS on the MAC domain can save.
The naive solution is to cost tiles without the allocation.
A planner that costs every tile at the worst-case precision~\cite{xue2026finedvfs}, P8 here, never underestimates a tile, but it overestimates every P4--P6 prefill block by up to $1.85\times$, treats TTFT slack as spent and finds no free points in compute-bound blocks.
Costing every tile at the allocation's average precision misjudges the blocks whose own precision differs from the average, though it ranks most steps correctly (\Cref{sec:results}).

